\documentclass[10pt,twocolumn]{article}
\usepackage[letterpaper,margin=0.75in,columnsep=0.25in]{geometry}
\usepackage{booktabs,amsmath,graphicx,microtype,hyperref,array}
\hypersetup{colorlinks=true,allcolors=blue}
\title{Checking or Recomputing? Two Cost Regimes in\\Verifying AI-Generated Formal Artifacts}
\author{Anonymous research draft}
\date{}
\begin{document}
\maketitle
\begin{abstract}
A formal proof can be cheap to check after an expensive search, yet a proof
that delegates its justification to computation can make the kernel repeat
that computation. We study this distinction in a released AI-generated Lean
library and use it to guide ProofServe, a tool for independently rechecking
compiled Lean projects. A probability sample of the pinned OpenAI mathematics
library currently contains 41 of 42 original build costs. The 20 complete
strata give a conditional compiler-work estimate of 278.09 CPU-hours, with a
fragile nominal 95\% interval of 94.15--462.04 hours; the global cost split
remains uncertain. Selected complete mathematical closures require 6.31\%
and 8.95\% of their compiler CPU for checking, whereas computation embedded in
proofs remains work for the checker. ProofServe combines original-kernel
declaration checking, dependency-aware parallelism, restartable segments and
machine-readable coverage and axiom receipts. Its mechanisms have precedents;
the systems contribution is their usable integration and evaluation.
This draft separates established measurement evidence from the pending
release-tool evaluation, which will determine its speed and external-project
acceptance.
\end{abstract}

\section{Introduction}
Lean's kernel validates elaborated proof terms rather than repeating the
tactic search that produced them. This separation motivates independent
checking of AI-generated artifacts. It does not make every check cheap:
definitional equality can evaluate a program embedded in a proof. At the
project boundary, a user must also establish which modules were checked,
which imports were trusted, and which axioms the accepted statements use.

The Lean FRO calls for routine external checking. This is already more than
a roadmap: Lean 4.35 ships \texttt{lake check}, which exports and rechecks
default targets, and can invoke bundled independent kernels
\cite{leanFroY4,lean435Check}. ProofServe addresses the operational cost and
scope of such a workflow, rather than claiming to invent independent
checking or to replace an independently implemented kernel.
Its intended user supplies a Lake project and receives both a checker verdict
and an explanation of its coverage.

Our empirical contribution is a pinned, reusable measurement package that
distinguishes source compilation, independent replay and computation.
Our systems contribution is ProofServe's project-facing check and audit
interfaces, together with a fixed evaluation of the release-form tool.
The latter is in progress; prototype speedups do not establish release
performance. The paper retains four questions:
\begin{description}
\item[RQ1] Where do verification costs fall, and how prevalent are the
computation-replaying and mathematical-proof-checking regimes across the library?
\item[RQ2] Why does independent verification fail to realize the potential
low cost of mathematical-proof checking?
\item[RQ3] Can computational checking scale across cores or a cluster
without expanding the trusted base?
\item[RQ4] What gaps remain between a reported verified result and build
coverage, reproducible outcomes and known axiom dependencies?
\end{description}

\section{Method and workload}\label{sec:method}
The primary workload is the published OpenAI mathematics repository
\cite{openaiMathSource}, pinned at \texttt{4eaee875}; the data package records
the full revision. It contains 163,097 source modules and uses Lean 4.34.1
with Mathlib \texttt{d13f23b7}. The sample has two original probability draws
in each of 21 strata, combining source-size percentile, source-import depth
and an explicit kernel-decision marker. The seed and all original identities
are fixed. Prerequisite modules and purposefully selected expensive examples
do not acquire probability-sample weights.

Source cost is native per-module compilation, including imports, tactics,
kernel checking and artifact emission; prerequisite preparation, queueing,
separate C compilation and Lake driver work are excluded. Checker clocks
include their declared import, restore and export lifecycle. CPU is user
plus system time, reported separately from wall time. Neither summed native
wall intervals nor CPU divided by worker count is project-build makespan.

Timed runs use the existing exclusive P-core lease on CPUs 0--7 of an
i9-14900K; untimed preparation uses E cores. Current Lean invocations omit
\texttt{-M}, including \texttt{-M0}. External charged-memory protection and
sampled RSS are reported separately. Our earlier use of Lean's nondefault
in-process memory checking and our prototype's missing persistence marking
were measurement-method defects, not inherent workload properties.
Host-reboot-interrupted attempts are retained but excluded from cost,
censoring and completed samples. Completed records are integrity-checked before reuse. Repeat order and
retained caches limit causal and cross-machine interpretation.

All measurement values below come from
\texttt{data/openai-math-4eaee875/}. Its small tables retain versions,
sample membership, missingness, clocks and estimator scope. Historical
prototype interventions appear only when explaining a cause.
Release system results must instead identify a committed executable under
\texttt{checker/results/}; external-project results belong under
\texttt{data/external-projects/}. We never substitute one evidence class for
the other.

\section{Two costs, uncertain population weights}\label{sec:costs}
\subsection{A proof term can contain substantial computation}
Ordinary mathematical proofs often store the result of expensive proof
search as a term whose validation is much cheaper. Reflection and
kernel-evaluated decision procedures can instead encode a calculation that
the kernel must perform to justify the proposition
\cite{lean4,leanDecide434}. A short source proof is therefore not a small
checking workload. Conversely, an expensive source tactic need not leave an
expensive term. This stage distinction, rather than authorship or source
length alone, defines the two cost regimes studied here.

The sample's static marker is any explicit comment/string-masked
\texttt{decide +kernel} occurrence. Its label is a reproducible source proxy,
not a complete semantic classifier: ordinary \texttt{decide}, definitional
reduction and generated tactics can also perform computation. Absence of the
marker does not establish inexpensive checking. A count of lexical sites
cannot be multiplied by a compiler-time fraction to obtain kernel work.

\subsection{Coverage does not imply precision}
At this snapshot, 41 of 42 original source costs have completed. One ordinary
draw remains unmeasured. Twenty complete strata cover 163,057 modules; the
remaining 40-module stratum has only one observed original draw and is
excluded from complete-cell inference. All nine marked strata are complete,
covering 5,848 modules. The extra fixed precision panels remain independent
of the original sample and do not enter the figures in this draft.

For complete strata, the stratified total is
\[
 \widehat T=\sum_h N_h\bar y_h,\qquad
 \widehat V=\sum_h N_h^2(1-n_h/N_h)s_h^2/n_h .
\]
The covered-domain compiler estimate is 278.09 CPU-hours: 57.55 in marked
modules and 220.55 in other modules. Thus the conditional marked share is
20.69\%, despite the much smaller marked-file population. This is compiler
work in marker-defined files, not the fraction of CPU spent on computation
inside proofs. It does not answer the full-library runtime-weight question.

The nominal 95\% Student-$t$ interval is 94.15--462.04 CPU-hours, using
Satterthwaite degrees of freedom 2.61. Two observations per cell and heavy
tails make the approximation fragile. It excludes the incomplete stratum,
nonresponse, runtime drift and transfer-model uncertainty. Its width is a
substantive result: nearly complete module coverage does not make a
two-draw-per-cell total precise.

A separate missing-cost scenario preserves all known original costs and
weights, assigning the last draw the observed ordinary-class minimum or
maximum. It gives 278.5--279.1 CPU-hours after outward rounding.
This narrow range concerns completion of the fixed sample, not its sampling
error. Without assumptions on the missing cost, the completion-cost upper
remains unbounded. Checker-transfer scenarios are models, not measured whole-library savings.

\paragraph{Design implication.}
An audit should expose the source proxy, measured coverage and interval
separately. It should never turn an unknown module into a zero-cost module or
a marker into a dynamic phase verdict. ProofServe therefore reports actual
checked scope and accepts costs only as explicitly supplied observations.

\section{Why cheap checking is hard to realize}\label{sec:coupling}
\subsection{Source production and independent replay have different work}
The complete selected closures in Table~\ref{tab:closures} show that checking
can be substantially cheaper than producing a mathematical artifact.
The checkpointed Falconer closure uses 2,038.44 CPU-seconds for checking
against 32,292.60 for compilation, or 6.31\%. A Mathlib closure under the
same checkpoint protocol uses 91.15 against 1,018.22 CPU-seconds, or 8.95\%.
Their wall ratios are 9.26\% and 14.67\%, respectively. CPU and wall answer
different questions and are not interchangeable speedups.

\begin{table}[t]
\centering
\small
\begin{tabular}{lrrr}
\toprule
Closure & Modules & CPU ratio & Wall ratio\\
\midrule
Mathlib Nat.Basic & 1,330 & 8.95\% & 14.67\%\\
Falconer & 13,205 & 6.31\% & 9.26\%\\
Grid & 3,793 & 8.47\% & 14.52\%\\
Amplitudes & 5,083 & 8.44\% & 14.55\%\\
\bottomrule
\end{tabular}
\caption{Measured complete-closure check/source-work ratios from the public
data package. Mathlib and Falconer use checkpointed replay; Grid and Amplitudes
use the older uninterrupted stream. Shared dependencies make these overlapping
cases, not independent population samples. Design implication: charge checking
preparation at the complete-closure boundary.}
\label{tab:closures}
\end{table}

Independent checking still requires producing its input artifact.
The source denominator includes imports, tactics and emission; replay
includes reading and its declared lifecycle. Their ratio is neither a
kernel fraction nor an interactive-service or project-rebuild speedup.

Whole-closure averaging also hides where the work sits. A closure includes
many reused mathematical dependencies around a computation-heavy target.
Grid and Amplitudes therefore do not establish that every expensive
certificate is cheap to recheck. Their aggregate CPU ratios cannot be
assigned to all marked modules as a measured population law.
This explains why we retain selected-case ratios while refusing a global
checking-cost point estimate.

\subsection{Imports and repeated preparation consume the opportunity}
Starting a checker independently for every module repeats process and import
work. Importing an entire project's final environment once avoids some
repetition, but it changes what is being checked. A module-local path may
trust its imports; a final merged environment may hide distinctions between
incoming definitions. Checking every module in its own declared dependency
context is a different ledger from checking one final map of constants.

The bundled Lean 4.34.1 checker already dispatches module-local tasks
\cite{lean434CheckerSource2026}. Its default path loads imports and replays
local declarations; \texttt{--fresh} replays the final imported environment
from an empty kernel environment. In the pinned implementation the latter
replay traverses that map sequentially. This observation is version-specific,
not a claim that official checking has no parallelism. Target-name equality
alone does not justify a performance ratio between these procedures and a
dependency-ordered per-module verifier.

These boundaries motivate a reusable checked prefix, but such reuse has a
trust cost. The reader, module visibility, duplicate handling and restored
metadata must preserve the intended context. Unchanged mathematical kernel
rules do not make the entire checker implementation unchanged.
ProofServe reports its execution path and actual accepted modules so a
consumer can distinguish these cases.

\subsection{Prototype interventions identify engineering causes}
Our prototype initially shared frontiers without explicit persistent marking.
The corrected same-executable intervention reduced eight-worker CPU by
14.13\% and wall time by 5.13\% on its measured closure. This corrects our own
implementation method. It does not establish that reference counting
explains every residual cost or that Lean intrinsically scales poorly.

A second prototype policy replayed an entire module serially when any
nonordinary declaration prevented its ordinary-declaration plan.
Retaining a dependency-closed nonordinary prefix and parallelizing the
remainder reduced the paired plain-run wall median from 257.66 to
196.67 seconds, or 23.67\%, at a 0.106\% CPU-median increase.
The prefix is a conservative implementation choice, not a proved serial
lower bound. The intervention's flags, version, full clocks and trust
boundary are recorded in the evidence map. These causal observations do
not fill the release-tool evaluation below, and their ratios are not
multiplied with other experiments.

\paragraph{Design implication.}
Independent checking should amortize compatible preparation, retain module
context and measure complete operations. Parallelism must come from actual
declaration dependencies, while memory lifetime and fallback policies remain
explicit implementation decisions. A faster leaf interval alone is
insufficient evidence of a faster project check.

\section{ProofServe design}\label{sec:design}
ProofServe exposes \texttt{check}, \texttt{audit} and the existing
Kimina-compatible \texttt{serve} interface. The service is maintained for
compatibility; this work adds no serving feature. Figure~\ref{fig:flow}
shows the project workflow. The native checker and audit implementation are
available; packaging and release-level acceptance remain in progress.

\begin{figure}[t]
\centering
\fbox{\begin{minipage}{0.91\columnwidth}
\centering
Lake project, toolchain and compiled imports\\[3pt]
$\downarrow$\\
\texttt{check}: inventory $\to$ dependency order\\
$\to$ original-kernel prefix and parallel waves\\
$\to$ checkpointed segments\\[3pt]
$\downarrow$\\
Receipt: scope, verdicts, axioms, CPU and wall\\[3pt]
$\downarrow$\\
\texttt{audit}: source/default targets + receipt\\
$\to$ omissions, lexical points, supplied costs
\end{minipage}}
\caption{Separate proof acceptance from its coverage explanation. Audit
consumes reported evidence; it does not independently certify the checker.}
\label{fig:flow}
\end{figure}

\subsection{Check a dependency closure}
\texttt{proofserve check <project>} uses the project's declared toolchain and
compiled artifacts. By default the native checker discovers compiled local
roots; explicit roots support custom layouts. It traverses transitive imports
and rejects missing artifacts instead of treating them as already checked.
Uncompiled source files remain outside the accepted compiled closure and
are exposed separately by audit.

For ordinary declarations, the checker constructs a local dependency graph
from their referenced constants. Ready declarations run against a checked
frontier using the original kernel's checking operation. Workers return
checked values; a coordinator validates their identities and commits the
frontier after a wave joins. Nonordinary declarations and their ordinary
prerequisites use the original replay path before the remaining ordinary
graph runs. Cyclic or unsupported local structure retains a conservative
fallback. No temporary axiom skeleton or repaired proof replaces the input.

Persistent marking occurs only after earlier workers join and before new
ones start. The release tool exposes disabling persistence and split fallback
for the fixed ablations. These flags test the integrated implementation;
prototype intervention results do not determine the outcome.
Worker scaling is measured in both wall and CPU; it need not be monotonic.

\subsection{Resume without relabeling reused work}
The checker processes bounded segments and serializes its successful prefix.
A repeated invocation can reuse completed segments when the toolchain,
artifacts, Lake lock and checkpoint configuration match. A live invocation
holds its checkpoint directory lock. An interruption does not create an
accepted segment; an incomplete receipt identifies only the successful
prefix. Restoration of an accumulated prefix still consumes memory, so
segmentation is not a guarantee that an arbitrarily large environment fits
a small memory budget.

Receipt v1 records the execution path, Lean and dependency identity, module
set hash, accepted modules, rejection events, axiom set and clocks.
Current-invocation CPU and reused cumulative segment CPU must be distinguished.
A completed checkpoint cache hit is not a fresh timed verification run.
The interruption experiment charges restart and restoration overhead.

\subsection{Audit what verified covers}
\texttt{proofserve audit <project>} inventories local source, extracts imports
and masks comments and strings before counting kernel-decision forms.
Supported declarative Lake configuration yields a default-target source
closure and omissions. Executable configuration, custom targets or
unsupported constructs produce an explicit unknown result unless the caller
supplies evaluated configuration or roots. Supplied roots are never relabeled
as Lake's defaults.

Audit can consume receipt v1 to display the checked compiled closure,
verdicts and reported axioms. An incomplete receipt establishes only the
axioms of its last successful prefix. Per-module cost logs require explicit
source digests; aggregate checker clocks are not redistributed into
invented module costs. Missing costs stay missing. Lexical classes and
measured costs remain separate fields, because neither source markers nor
an accepted theorem establish its dynamic regime.

The current receipt does not bind its compiled evidence to the current
source digest. Audit makes that gap explicit rather than certifying a
changed source tree using an old receipt. Generated modules, external source trees and executable Lake behavior
also require separate coverage evidence.

\subsection{Trust and compatibility}
ProofServe independently reruns a workflow using Lean's original kernel;
it is not a separately proved or independently implemented kernel.
It additionally trusts the artifact reader, dependency/context metadata,
checked-token FFI adapter, quiescent persistent marking and its own locally
generated checkpoints. Serialized private metadata and an internal Lean
symbol require a matching SDK. Input hashes detect ordinary staleness;
they do not authenticate maliciously modified checkpoints.

Axioms remain axioms. Acceptance does not mean axiom-free, absence of
\texttt{sorryAx}, correct tactic plugins, or correct executable code.
Upstream unsafe/partial exclusions remain visible. On an unsupported or
mismatched toolchain, the command must fail explicitly or select the
project's matching official checker, with that path and its narrower
receipt scope identified. It must never silently decode another version's
artifacts or count official fallback as a parallel-tool measurement.

\section{Release-tool evaluation}\label{sec:evaluation}
The fixed evaluation answers the systems portions of RQ3 and RQ4. It is
limited to the committed \texttt{proofserve-check} executable and
\texttt{proofserve audit}. Every system row must identify the executable
commit, project commit, toolchain, roots, accepted module set and timing
boundary. Table~\ref{tab:release} reserves the result layout. The available
one-point record predates the frozen final matrix and is excluded.
No release speedup is established in this draft.

\subsection{RQ3: complete scaling, ablations and official scope}
Mathlib Nat.Basic, Grid and Barrier are checked at 1, 2, 4 and 8 workers,
with three fresh-checkpoint processes per point. Mathlib in that matrix
means the Nat.Basic closure, not the whole source tree; any separate
whole-Mathlib run has its own scope and row. We report wall and CPU medians,
ranges and matched acceptance sets. Each cell includes process startup,
artifact reading and its checkpoint lifecycle. Queueing for the common
P-core lease is outside the native clock and recorded separately.

\begin{table}[t]
\centering
\small
\begin{tabular}{lccc}
\toprule
Closure & Workers & CPU / wall & Set\\
\midrule
Mathlib Nat.Basic & Pending & Pending & Pending\\
Grid & Pending & Pending & Pending\\
Barrier & Pending & Pending & Pending\\
\bottomrule
\end{tabular}
\caption{Reserved release-matrix layout: three fresh-checkpoint repeats
per worker count. Cells await the frozen executable's complete results;
historical prototype values are not substituted.}
\label{tab:release}
\end{table}

At eight workers, Grid and Mathlib each receive three runs without
persistence and three without split fallback. Other settings and inputs match. These ablations test whether prototype
causes survive integration; CPU and memory accompany wall improvements.

Official \texttt{leanchecker} default and \texttt{--fresh} runs use the same
resource point and complete recorded command clocks. Default name-prefix
selection can include extra modules, and fresh mode can have a different
declaration ledger. We compare acceptance only at a demonstrated common
scope and report any unmatched work; where equivalence is absent, both
costs are shown without a speedup ratio. Earlier leaf-only official timings
are not an admissible full-closure denominator. The fixed comparison does
not establish superiority to every current official export or independent
kernel path.

RQ3's multi-core result retains additional trusted coordination code;
it cannot establish an unchanged implementation TCB. Cluster checking is
unevaluated and cannot be inferred from single-host curves.

\subsection{RQ4: acceptance, rejection, resumption and external use}
Development controls exercise loose bound variables, unknown constants and
wrong sorts, and compare small accepted closures with the official checker.
These are finite functional checks, not the completed release CI or a
general soundness proof. The release test requires all three bad-proof
classes to be rejected, accepted module sets and reported axioms to agree
at the comparable scope, and the interruption/resume run to match its
uninterrupted counterpart. Additional CPU and wall time from restart are
reported as costs rather than omitted from the successful run.

\begin{table}[t]
\centering
\small
\begin{tabular}{>{\raggedright\arraybackslash}p{0.48\columnwidth}>{\raggedright\arraybackslash}p{0.39\columnwidth}}
\toprule
Release acceptance evidence & Current disposition\\
\midrule
Three bad-proof classes in CI & CI closure pending\\
Interrupted vs.\ uninterrupted & Timed pair pending\\
Three unseen public projects & Results pending\\
At least one published AI project & Provenance pending\\
Package/container, no Workspace paths & Clean-start acceptance pending\\
\bottomrule
\end{tabular}
\caption{External and functional acceptance obligations. Pending entries
are not zero costs, failed proofs or completed user acceptance.}
\label{tab:acceptance}
\end{table}

External validation uses at least three public Lean projects not used during
development, including at least one with published AI-generated source.
For each, the result records its name and immutable commit, AI provenance,
module count, tool and official CPU/wall costs, acceptance parity, audit
findings and encountered problems. A checker fixture alone does not establish AI provenance. Clean-start
package/container execution must work without Workspace paths.

Audit acceptance examines both known and unknown coverage: default roots
that omit local modules, unsupported executable configuration, source
changes after a receipt, incomplete-prefix axiom reporting and supplied
cost-log mismatch. It must preserve explicit zero counts while distinguishing
an incomplete scan from an empty finding. These are user-visible coverage findings for RQ4.

\paragraph{Artifact and reproduction.}
The artifact is at \url{https://github.com/yunwei37/proofserve}.
The intended installed interface is
\texttt{proofserve check /path/to/project} followed by
\texttt{proofserve audit /path/to/project}.
The current native entry point is \texttt{proofserve-check}, built by
\texttt{cd checker; lake build}; audit also has the portable entry point
\texttt{python3 tools/audit/audit.py}.
The data-only command
\texttt{python3} \path{data/openai-math-4eaee875/summarize.py}
recomputes the published sample summary without Lean or raw Workspace
captures. Release installation, complete CI and external results remain
acceptance requirements, not claims made by these commands.

\section{Limitations and related work}
The current sample is incomplete, its static labels are proxies, and its
nominal sampling interval is wide. Closure ratios cover selected,
overlapping cases and retain their original protocol differences.
They do not identify global kernel phase weights or the checking cost of
the entire AI library. Release evaluation and external acceptance are
unfinished; their reserved cells will be filled only from the designated
committed tools. Hardware, cache order, observers and source/compiled
binding limit transfer beyond the measured environment.

Parallel proof checking, asynchronous tasks and certified decomposition
have substantial precedent in Isabelle, Coq and Lean
\cite{wenzelParallel2009,coqAsync2015,leanParallelRelease2025}.
Existing Lean checkers and the shipped Lake workflow address trust and
rechecking \cite{leanCheckerDocs,lean435Check}; ProofServe does not claim
mechanism firstness. Kimina and AXLE already provide verification services
\cite{kiminaPaper,axle}, so a new serving feature is not a contribution here.
A public technical report also studies build/check costs and computation
in the same OpenAI library \cite{openaiMathTooling2026}; this paper claims
neither the first same-corpus characterization nor the first observation
that checking differs from recomputation.

\section{Conclusion}
Proof checking and proof-embedded recomputation create different
opportunities, while preparation and coverage determine whether an apparent
saving survives at the project boundary. The pinned data package establishes
selected cheap-checking cases and a population estimate whose uncertainty
remains large. ProofServe turns those lessons into an explicit check/audit
workflow with original-kernel checking and declared additional trust.
Release scaling, correctness and external acceptance remain the test of
its practical value.

\bibliographystyle{plain}
\bibliography{references}
\end{document}
